\chapter*{Introduction to deduction}
\addcontentsline{toc}{chapter}{Introduction to deduction}
\markboth{Introduction to deduction}{Introduction to deduction}
\label{intro:deduction}

A deductive science does not look at the world. It fixes a small number of
starting rules, the \emph{axioms}, and constructs everything else from them.
Whatever it proves is true \emph{relative to those axioms}, without any
observation. This is the sense in which a deductive science constructs its world.

\section*{The one rule}

\begin{center}
  \emph{Nothing may be used before it has been constructed.}
\end{center}

\noindent Every chapter uses only definitions and results from earlier chapters, or
from earlier in the same chapter. When a later topic needs something new, the new
thing gets its own definition first. So this part of the book is one single
chain that can be read, and checked, from the beginning.

\section*{Where the chain starts}

Something has to come first, and it can't be constructed from anything earlier.
We start with:
\begin{enumerate}
  \item \textbf{Classical logic}, the rules of correct reasoning
    (\cref{ch:logic}). It is the language everything else is written in.
  \item \textbf{Zermelo--Fraenkel set theory}: one kind of object, the
    \emph{set}, and one relation between objects, \emph{membership} $\in$
    (\cref{ch:sets}).
\end{enumerate}
Everything after that is constructed: pairs, relations and maps, then the number
systems
\[
  \N \subset \Z \subset \Q \subset \R,
\]
then sequences, functions, calculus, linear algebra, and from there the rest of
mathematics.

\section*{Discrete mathematics first}

Discrete mathematics deals with objects that are separate and can be counted:
truth values, finite collections, whole numbers, fractions. It comes first
because the rest of mathematics is built on it. The real numbers, and with them
functions, series and calculus, are constructed from the rational numbers, and
those are constructed in the first chapters:
\[
\begin{array}{ccccccccc}
  \text{logic} & \to & \text{sets} & \to & \text{pairs} & \to &
  \text{relations} & \to & \text{maps} \\[1ex]
  & & & & & & \downarrow & & \downarrow \\[1ex]
  \Q & \leftarrow & \Z & \leftarrow & \N & & \text{equivalence classes}
  & & \text{recursion}
\end{array}
\]
\begin{itemize}
  \item A \textbf{natural number} is a particular set: $0 = \varnothing$, $1 =
    \set{0}$, $2 = \set{0, 1}$, and so on. Each number is the set of all numbers
    before it (\cref{ch:naturals}).
  \item An \textbf{integer} is what a difference $a - b$ of natural numbers
    \emph{would} be. Subtraction doesn't exist yet, so we use the pair $(a, b)$
    and treat two pairs as equal when they would give the same difference
    (\cref{ch:integers}).
  \item A \textbf{rational number} is built the same way from a quotient $a / b$
    of integers (\cref{ch:rationals}).
\end{itemize}

\noindent Planned chapters, each using only what is already in place:
\begin{itemize}
  \item finite and infinite sets, cardinality, Cantor's theorem;
  \item combinatorics: sums, products, binomial coefficients;
  \item divisibility and primes: Euclid's algorithm, positional notation, unique
    factorization;
  \item graphs;
  \item the real numbers, sequences and series;
  \item functions and calculus;
  \item linear algebra.
\end{itemize}

\section*{Conventions}

\begin{itemize}
  \item Numbered statements are counted per section: Definition~3.2.4 is the
    fourth numbered item of Section~3.2. Axioms, definitions, theorems, lemmas,
    corollaries and remarks all share one counter, so references are easy to
    find.
  \item A \emph{theorem} is an important result, a \emph{lemma} a stepping
    stone, a \emph{corollary} a direct consequence.
  \item A reference such as \enquote{Lemma~5.2.2(3)} means item~(3) of Lemma~5.2.2.
  \item Every proof ends with $\blacksquare$.
  \item Notation follows the \hyperref[rules]{Rules of notation}: one notation
    for one thing, and every notation is listed in the
    \hyperref[register]{Register}.
  \item Exercises are optional. Nothing later depends on them.
  \item Examples and remarks may mention things that come later, to motivate or
    illustrate. Definitions and proofs never use them.
\end{itemize}
